% VARIANT — not \input by main.tex; promote manually
% rvv-strip-mining_twotone.tex — alternative colour strategy: a clean two-hue
%   alternation between consecutive iterations (emphasises iteration BOUNDARIES
%   rather than per-iteration identity). Remainder still tail-tinted. Tokens only.
\tikzsetnextfilename{ch3-riscv-rvv-strip-mining_twotone}
\begin{tikzpicture}[
  cell/.style={draw=figrule, figLW/hair, minimum width=0.45cm, minimum height=0.5cm,
               font=\ttfamily\tiny, inner sep=0pt},
  lbl/.style={figT/label},
  micro/.style={figT/annot},
  count/.style={figT/annot, anchor=west},
  mbrace/.style={decorate, figLW/base, draw=figrule,
                 decoration={brace, amplitude=4pt, mirror, raise=1pt}},
  even/.style={fill=figprocess!24}, odd/.style={fill=figaccent!24},
  tail/.style={fill=figthreat!22},
]
% === VLMAX = 4 (top) ===
\node[lbl, anchor=west] at (-1.6, 0.8) {\textbf{VLMAX = 4}\ \ (\VLEN=256, SEW=64, LMUL=1)};
\foreach \i in {0,...,19} {
  \pgfmathsetmacro{\xpos}{\i * 0.5}
  \pgfmathtruncatemacro{\grp}{int(\i/4)}
  \ifodd\grp \node[cell, odd] at (\xpos, 0) {\i};
  \else \node[cell, even] at (\xpos, 0) {\i};\fi
}
\foreach \start in {0,2,4,6,8} {
  \draw[mbrace] (\start-0.2, -0.4) -- (\start+1.6, -0.4)
    node[midway, below=3pt, micro] {vl=4};
}
\node[count] at (10.0, 0) {5 iterations};

% === VLMAX = 8 (bottom) ===
\node[lbl, anchor=west] at (-1.6, -2.0) {\textbf{VLMAX = 8}\ \ (\VLEN=512, SEW=64, LMUL=1)};
\foreach \i in {0,...,19} {
  \pgfmathsetmacro{\xpos}{\i * 0.5}
  \ifnum\i>15
    \node[cell, tail] at (\xpos, -2.8) {\i};
  \else
    \pgfmathtruncatemacro{\grp}{int(\i/8)}
    \ifodd\grp \node[cell, odd] at (\xpos, -2.8) {\i};
    \else \node[cell, even] at (\xpos, -2.8) {\i};\fi
  \fi
}
\draw[mbrace] (-0.2, -3.2) -- (3.6, -3.2) node[midway, below=3pt, micro] {vl=8};
\draw[mbrace] (3.8, -3.2) -- (7.6, -3.2) node[midway, below=3pt, micro] {vl=8};
\draw[mbrace, draw=figthreat, figLW/emph] (7.8, -3.2) -- (9.6, -3.2)
  node[midway, below=3pt, micro, text=figthreat!80!black] {vl=4 (remainder)};
\node[count] at (10.0, -2.8) {3 iterations};
\end{tikzpicture}
